\section{CZARY}

Ogólna zasada:

Postacie, posiadające potencjał magiczny większy od zera mogą rzucać czary.
Potencjał magiczny każdej postaci składa się z trzech części, a każda z tych części jest używana wtedy, gdy dominującym jest odpowiadające jej słońce.
Dominujące słońce jest ustalane według procedury, opisanej w punkcie 4.0.
Wpływa ono na liczbę czarów, których dana postać może używać podczas jednej przygody.
Rodzaj czarów decyduje o czasie, w którym może on być użyty.

Sposób postępowania:

Przed rozpoczęciem gry każdy gracz musi wybrać dla swej postaci liczbę czarów równą najwyższej z trzech wartości, składających się na potencjał magiczny tej postaci oraz zapisać nazwy tych czarów na jej karcie cech.
Powstaje w ten sposób pula czarów, z której wybierane są te, które będą używane podczas bieżącej przygody (patrz 4.1 D).
Postacie rzucając czary, odnoszą przy tym rany w liczbie równej kosztowi danego czaru.
Pewne czary mogą zostać odparte, a tym samym ich efekt zostaje zniwelowany (patrz 11.0).

Zasady szczegółowe:

[10.1] Czary zostały podzielona na typy, w zależności od sytuacji, w której są najbardziej efektywne.
Wszystkie czary, poza typem W, mogą być rzucane w dowolnym momencie podczas gry.
Typ W może być używany tylko podczas walki. (Patrz tabela 10.9).

[10.2] Postać może używać danego czaru dowolną liczbę razy, pod warunkiem, że posiada jeszcze dostateczną siłę, by koszt rzucenia danego czaru jej nie zabił.

[10.3] Postać, używająca w czasie walki czarów nie może podczas tej samej rundy atakować przy pomocy broni.

[10.4] Podczas walki czary mogą być rzucane tylko przez postacie, znajdujące się w drugim szeregu szyku oddziału.
Potwory mogą rzucać czary zarówno z drugiego, jak i z pierwszego szeregu swego szyku.

[10.5] Efekty poszczególnych czarów typu W (używanych tylko w czasie walki):

Urok (koszt: 3) – Może być rzucony na każdego potwora, znajdującego się w pierwszym lub drugim szeregu szyku potworów.
Jeżeli potwór nie odeprze uroku natychmiast traci wolną wolę i przechodzi pod całkowitą kontrolę postaci, która ten czar rzuciła.
Taki potwór staje się, zarówno pod względem intencji, jak i rzeczywistego zachowania, członkiem oddziału.
Powinien być natychmiast umieszczony w wolnym miejscu w szyku oddziału.
Jeżeli jednak rzucała urok zostaje zabita czar traci moc i potwór natychmiast atakuje pozostałych członków oddziału – jeżeli dzieje się to podczas walki należy go umieścić w szyku potworów, jak najbliżej oddziału.
Potwór, znajdujący się pod wpływem uroku, zabijając inne potwory nie zdobywa punktów doświadczenia – punkty te przypadają należącym do oddziału postaciom.
Uwaga – wszystkie powyższe zasady (z pewnymi wyjątkami – patrz 15.5) znajdują zastosowanie w sytuacji, w której pod wpływem uroku, rzuconego przez potwora znajdzie się członek oddziału.

Wybuch (koszt: 1) – Może być rzucony na potwora, znajdującego się w pierwszym szeregu szyku potworów.
Jeżeli potwór nie odeprze tego czaru odnosi dwie rany.

Eksplozja (koszt: 1) – Działa na wszystkich, którzy są zaangażowani w walkę.
Wszystkie postacie i potwory, które go nie odeprą odnoszą 1 ranę.

Błyskawica (koszt: 2) – Może być rzucony na potwora w pierwszym szeregu szyku i jeżeli nie zostanie odparty zadaje 1K3x2 rany.

Uśpienie (koszt: 2) – Może być rzucony na każdego potwora w każdym szeregu i jeżeli nie zostanie odparty powoduje natychmiastowe zapadnięcie tego potwora w sen.
Uśpiony potwór może zostać zabity bez potrzeby rzucania kostką (oczywiście po pozbyciu się w ten czy inny sposób wszystkich pozostałych potworów).

Uwolnienie (koszt: 3) – Pozwala postaci, na którą został rzucony urok i która teraz, podczas walki, występuje przeciwko oddziałowi na powrót w jego szeregi czyli niweluje działanie uroku.
Postać pod wpływem uroku musi próbować odeprzeć ten czar.

Magiczna tarcza (koszt: 2) – Rezultatem tego czaru jest całkowita odporność postaci, na którą został rzucony na wszystkie czary, zadające rany na okres jednej walki.

Niepewność (koszt: 2) – Daje oddziałowi czas na użycie tuż przed rozpoczęciem walki mikstur i czarów uzdrawiających.
Nie może zostać odparty.

Zawieszenie broni (koszt: 3) – Jego skutek jest analogiczny do porozumienia w czasie negocjacji (patrz 8.6), czyli czar ten powoduje natychmiastowe przerwanie walki.
Może być użyty w czasie każdej fazy walki i nie może zostać odparty.
Nie wolno go stosować w czasie starcia z Bezimiennym.

Psychiczny cios (koszt: 4) – Może być rzucony na potwora w pierwszym szeregu szyku i jeżeli nie zostanie odparty zadaje 2K6+2 ran.

[10.6] Czary innych, poza W, typów mogą być rzucane w dowolnym momencie podczas gry.
Pewne z nich powinny być jednak używane w określonych sytuacjach, gdyż poza nimi ich efekt jest zerowy.

Zamek (koszt: 1) – Może być użyty do zablokowania drzwi.
Rzucany jest zwykle po otwarciu drzwi i po stwierdzeniu, że w segmencie, do którego oddział miał zamiar wejść znajduje się potwór.
Zamiast negocjować, wykupywać się lub toczyć walkę oddział może się wycofać i przy pomocy tego czaru zamknąć za sobą drzwi.
Takie drzwi już nigdy nie będą mogły być otworzone.
Jeżeli oddział wejdzie do tego samego segmentu później i z innej strony musi tam spotkać tego samego potwora.

Magiczna zbroja (koszt: 1) – Ten czar tworzy pole ochronne wokół postaci, na którą został rzucony.
Należy rzucić 1K3+1 i zapisać wynik.
Za każdym razem, gdy postać, znajdująca się pod ochroną tego czaru odniesie ranę ten wynik należy redukować o jeden.
Gdy osiągnie on zero wszystkie następne rany postać odnosi już normalnie.
Magiczna zbroja nie chroni przed ranami, będącymi kosztem rzucania czarów, ani nie chroni przed czarami nie zadającymi ran (takimi jak urok).

Neutralizacja trucizny (koszt: 1) – Ten czar może być rzucony na postać, która wypiła truciznę, bada zatrutą fontannę lub wpadła w pułapkę, zawierającą truciznę.
Neutralizuje wszelkie ujemne skutki działania trucizny.
Musi być rzucony natychmiast po zatruciu.

Ożywienie kamienia (koszt: 3) – Ten czar może być użyty w celu przywrócenia do życia postaci, która została zamieniona w kamień przez Meduzę.
Nie może być rzucany podczas walki.
Ożywiona postać posiada tę samą liczbę ran, jaką miała przed zamienieniem jej w kamień.
Nie może być rzucony na samego siebie.

Siła (koszt: 1) – Ten czar ma taki sam efekt, jak mikstura ,,Zręczność'' – patrz 14.5

Teleportacja (koszt: 3) – Ten czar teleportuje dowolną postać lub potwora do wskazanego przez rzucającego, odwiedzonego wcześniej przez oddział segmentu labiryntu (potwór zostaje w tym segmencie).
Nie może być użyty podczas walki.
Nie może być użyty przeciwko Bezimiennemu, Demonom oraz Czarnym Wrotom.

Uzdrowienie (koszt: 1) – Leczy postać, na którą jest rzucony z 1K3+1 ran.
Nie może być rzucony na siebie samego.
Uwaga! Zanim postać, która może rzucać ten czar wejdzie do labiryntu, należy, przy pomocy rzutu 1K6, ustalić ile razy ten czar będzie mógł zostać użyty podczas bieżącej przygody.
Po wyczerpaniu limitu, określonego liczbą wyrzuconych oczek postać nie może już rzucać czaru uzdrowienia (jest to wyjątek od zasady 10.2).

Odmłodzenia (koszt: 2) – Ten czar ma taki sam efekt jak uzdrowienie, z tą różnicą, iż leczy z 1K6+1 ran.
Również może być użyty tylko ograniczoną liczbę razy w czasie jednej przygody. (Ograniczenie według zasady jak przy uzdrowieniu).

Delikatne Palce (koszt: 1) – Ten czar zwiększa zdolności ,,pułapki'' o 3 punkty, ale tylko na czas, potrzebny na przeprowadzenie jednej próby unieszkodliwienia pułapki.
Nie może być rzucony na postać, która w ogóle nie posiada zdolności ,,pułapki''.

[10.7] Czary, pomagające w negocjacjach (typ N) mogą być używane przed rzutem kostką, przesądzającym o wyniku negocjacji, zaś czary typu V – podczas próby wykupienia się.

Elokwencja (koszt: 1) – Ten czar pozwala na dodanie 4 punktów do liczby oczek, wyrzuconej podczas negocjacji (dodatkowo, niezależnie od wszelkich modyfikacji, wynikających ze zdolności negocjującej postaci).

Zastraszenie (koszt: 2) – Daje efekt (bez konieczności rzucania kostką) analogiczny do ,,zastraszenia'' w negocjacjach (patrz 8.6).

Zniechęcenie (koszt: 3) – Ten czar daje podobny efekt, jak zastraszenie, z tym, że potwór oddaje wszystkie swoje skarby.

Sugestia (koszt: 1) – Pozwala oddziałowi na odjęcie 2 punktów od liczby oczek, wyrzuconych przy próbie wykupienia się.

Pochlebstwo (koszt: 2) – Pozwala oddziałowi na odjęcie 4 punktów od liczby oczek, wyrzuconych przy próbie wykupienia się.

[10.8] Umiejętność rzucania czarów specjalnych (typ S) może być zdobyta tylko w toku gry, tzn. nie mogą one być wybierane na początku, do ,,wstępnej puli'' czarów.

Gniew Boży (koszt: 3) – Ten czar może być użyty tylko podczas walki.
Umiejętność jego rzucania jest nabywana poprzez znalezienie manuskryptu lub ołtarza Malthusa (patrz 13.9).
Jeżeli nie zostanie odparty zadaje 2K6+2 ran.

Wskrzeszenie (koszt: 5) – Tego czaru nie wolno używać podczas walki.
Umiejętność jego rzucania może być nabyta tylko przez znalezienie pierścienia ,,Wskrzeszenie'' (patrz 14.9) lub regału z księgami (patrz 13.9).
Może być użyty do wskrzeszenia postaci, zabitej podczas walki.
Musi być rzucony natychmiast po zakończeniu walki, w której postać poniosła śmierć, gdyż użyty później nie będzie działał.
Wskrzeszona postać ma tyle ran, ile miała w momencie wejścia do labiryntu.

[10.9] Tabela ,,Zestawienie czarów'' – patrz plansza z tabelami.
